
\begin{proposition}[Every uncontracted degree-one horizontal image has nonnegative Euler number]\label{cdp:st:C.2.RAT.3}
Retain the torus-fibre, rational-normalization and ruled-boundary
morphism of \ref{cdp:st:C.2.RAT.1}. Suppose that $C$ is not contracted and
that its actual map to the normalization of
$H=h(C)_{\rm red}$ has degree one. Let $r_{\rm dom}$ be the
number of original ruling fibres $F_j$ whose positive vertical
boundary contains a strict conductor component mapping onto
$H$. Then
\[
       0\le r_{\rm dom}\le2,\qquad
       e(S)\ge2-r_{\rm dom}\ge0.
\tag{C.RA.21}\label{cdp:eq:C.RA.21}
\]
There is no assumption on the generic transverse type, on the
local complete intersection property of $H$, or on the
nilpotent structure of the conductor image.
\end{proposition}
\begin{proof}
The conductor coefficient on $C$ is two, because $C$ is
the strict transform of the coefficient-two section in
(C.RA.1) and is not $\pi$-exceptional. At a general point
of $H$, its degree-one map supplies one branch of the reduced
transverse plane curve. If that branch has intrinsic
conductor exponent $c_i^{\rm int}$, its exponent in the
whole plane curve is
\[
             2=c_i^{\rm int}+\sum_{j\ne i}I(C_i,C_j).
\tag{C.RA.22}\label{cdp:eq:C.RA.22}
\]
Here $I(C_i,C_j)$ is the positive intersection multiplicity
of two different plane branches. To verify the displayed
comparison, write the reduced equation as $u\prod f_j$.
On the normalization of the $i$th branch the residue
description of its dualizing generator gains the factor
$\prod_{j\ne i}f_j|_{C_i}$; its valuation is the sum in
\eqref{cdp:eq:C.RA.22}. The remaining factor is exactly the conductor
of the individual branch. This gives the actual inclusion
of its conductor ideal into its normalization, with the
stated exponent. Thus there are at most two other generic
branches. Each vertical strict component dominating $H$
supplies at least one such branch, counted with its covering
degree. This proves $r_{\rm dom}\le2$.

For an original fibre $F_j$, let $B_j$ be the positive
vertical boundary in its total transform. Its reduced
support is connected and meets $C$ at a single point.
These assertions follow at each point blowup: a centre
at the intersection of positive components inserts a
positive exceptional component; a centre at a smooth
coefficient-one point creates only a coefficient-zero
leaf; and a centre at a smooth point of coefficient
at least two creates a positive leaf. These are all the
possibilities in the exact divisor formula
$D_{\rm new}=\sigma^*D_{\rm old}-E$.
The $r$ points of $C$ so obtained are distinct.

If every component of $B_j$ is $\pi$-exceptional, there
has been at least one multiplicity-one blowup centre
above $F_j$. Indeed, in the absence of such a centre,
every component of the entire total transform of $F_j$
has positive conductor coefficient. That complete
ruling fibre, with its positive multiplicities, has
self-intersection zero. It cannot be supported on the
negative-definite exceptional matrix of a normal surface
resolution. Distinct original fibres require distinct
blowup events. Hence at most $N_0$ of the $B_j$ can be
entirely $\pi$-exceptional.

Discard those fibres and the $r_{\rm dom}$ fibres counted
above. There remain at least
\[
                        k\ge r-N_0-r_{\rm dom}
\tag{C.RA.23}\label{cdp:eq:C.RA.23}
\]
fibres with a surviving strict conductor component and
none dominating $H$. Follow the connected positive
boundary from its point on $C$ to the first component
which is not $\pi$-exceptional. Its image is a curve
different from $H$ meeting $H$ at that point's image.
Every intervening connected chain is contracted by
$\pi$, so this is an actual intersection in $S$, not
an identification imposed on an auxiliary graph.

Use the normalization-incidence graph of
$G=\Gamma_{\rm red}$: its vertices are the rational
normalization components and the singular points with
more than one normalization branch, and its edges are
the individual branch preimages. The preceding $k$
intersections give $k$ distinct edges adjacent to the
vertex belonging to $H$. They remain distinct even
when two of their images are the same singular point:
$C\to H^\nu$ is an isomorphism, and the points on $C$
were distinct. Remove that vertex and all its adjacent
edges. Each of the $c$ remaining connected components
which receives one of these edges contains a
normalization vertex of a curve other than $H$.
In particular $c\le n_G-1$, where $n_G$ is the number
of irreducible components of $G$.

Contract a spanning tree in each of those $c$ graph
components. The $k$ retained edges then give a connected
multistar with $c+1$ vertices and $k$ edges. Its first
Betti number is $k-c$, so the first Betti number
$\beta$ of the full incidence graph satisfies
\[
             k-c\le\beta\le1,
          \qquad n_G\ge k+1-\beta\ge k.
\tag{C.RA.24}\label{cdp:eq:C.RA.24}
\]
The upper bound on $\beta$ is the actual local torus-group
quotient (A.RT.18)--(A.RT.20). Adding genus-zero
normalization components changes no graph first Betti
number, and the normalization quotient gives
$e(G)=n_G+1-\beta\ge n_G$. Combining \eqref{cdp:eq:C.RA.23}--\eqref{cdp:eq:C.RA.24}
with the exact surface Euler formula yields
\[
\begin{split}
 e(S)&=2-r+N_0+e(G)\\
     &\ge2-r+N_0+n_G
      \ge2-r+N_0+(r-N_0-r_{\rm dom})
       =2-r_{\rm dom}.
\end{split}
\tag{C.RA.25}\label{cdp:eq:C.RA.25}
\]
This proves the assertion.
\end{proof}

The degree-one result accounts for every generic transverse
type allowed by conductor exponent two on that branch. It
does not replace the other branches by a presumed ordinary
triple: those branches may have larger conductor exponents.
The cases in which the horizontal section is contracted, or
maps with higher degree to a singular conductor image, require
the further arguments identified separately below.


